\section{Ambient adjacency access}
\label{sec:adjacency}

Consider an undirected graph with vertex labels.  A standard adjacency-list
interface may provide the following operations, with returned information and
work explicitly charged:

\begin{enumerate}[label=(\roman*)]
\item query the degree of a known vertex;
\item query one indexed neighbor and its edge weight, or scan the complete
      adjacency list at degree-proportional cost;
\item store arbitrary data derived from previous replies;
\item output a list of vertex labels and numerical values.
\end{enumerate}

The graph labels outside the revealed boundary carry no usable structure.  An
adaptive algorithm chooses its next query from the transcript.  Randomized
algorithms additionally have private random bits, and the lower-bound theorem
must state whether the graph or query adversary is oblivious or adaptive.

This access model says nothing about the arithmetic after a scan.  In
particular, after exposing an active principal matrix $Q_{SS}$, the algorithm
may legally:

\begin{itemize}
\item apply one row or one sparse matrix--vector product;
\item run local coordinate or proximal updates;
\item construct a Krylov basis;
\item factor $Q_{SS}$ or an associated Laplacian;
\item eliminate vertices and retain fill or response messages;
\item build a spectral sparsifier or preconditioner;
\item answer future queries from any persistent representation it has paid to
      construct.
\end{itemize}

Therefore ``adjacency access'' and ``first-order method'' are complementary
phrases, not competing definitions.  The first describes the input oracle; the
second restricts how numerical information is generated.
